
\documentclass{article}
\usepackage{colortbl}
\usepackage{makecell}
\usepackage{multirow}
\usepackage{supertabular}

\begin{document}

\newcounter{utterance}

\centering \large Interaction Transcript for game `codenames', experiment `risk\_low', episode 5 with reflexion{-}Qwen3.5{-}9B.
\vspace{24pt}

{ \footnotesize  \setcounter{utterance}{1}
\setlength{\tabcolsep}{0pt}
\begin{supertabular}{c@{$\;$}|p{.15\linewidth}@{}p{.15\linewidth}p{.15\linewidth}p{.15\linewidth}p{.15\linewidth}p{.15\linewidth}}
    \# & $\;$A & \multicolumn{4}{c}{Game Master} & $\;\:$B\\
    \hline

    \theutterance \stepcounter{utterance}  
    & & & \multicolumn{2}{p{0.3\linewidth}}{
        \cellcolor[rgb]{0.9,0.9,0.9}{
            \makecell[{{p{\linewidth}}}]{
                \texttt{\tiny{[GM$|$GM]}}
                {'hidden': {'team': ['teacher', 'phoenix', 'plate'], 'innocent': ['heart', 'plastic', 'war', 'ninja', 'millionaire', 'snowman', 'orange'], 'opponent': ['root'], 'assassin': []}, 'revealed': {'team': {'team': [], 'innocent': [], 'opponent': [], 'assassin': []}, 'opponent': {'team': [], 'innocent': [], 'opponent': [], 'assassin': []}}}
            }
        }
    }
    & & \\ \\

    \theutterance \stepcounter{utterance}  
    & & \multicolumn{4}{p{0.6\linewidth}}{
        \cellcolor[rgb]{0.9,0.9,0.9}{
            \makecell[{{p{\linewidth}}}]{
                \texttt{\tiny{[P1$\langle$GM]}}
                \texttt{You have attempted this same task before and did not succeed. These are the plans you wrote after those attempts.} \\
\texttt{Trial \#1: New plan: I failed because I included "phoenix" in the TARGETS list for the clue "school," which is semantically incorrect and likely violated the rule that the clue must be related to the target words. Additionally, I did not verify that the clue word itself was not a distractor or opponent word before submitting. Next time, I will strictly select a clue word that has a clear, direct semantic link to all listed targets (e.g., "education" for "teacher") and double{-}check that the clue word is not present in the opponent's list, distractor list, or kill words before formatting the response. I will also ensure the TARGETS list only contains words from my remaining team list that are actually related to the clue.} \\
\\ 
\\ 
\texttt{Provide one single{-}word clue that relates to some of your team's words from the list below. You can choose to how many and to which words your clue relates to.} \\
\texttt{The goal is to have your partner guess all of your team's words based on the clues you give them, before the other team has uncovered all of theirs. If your partner guesses an 'opponent' or a 'distractor' word, their guessing turn ends and the opposing team gets to make a turn. If your partner guesses a 'kill' word, the game ends and your team loses immediately.} \\
\texttt{The more words your clue relates to, the quicker you can win the game, but the harder the associations may be to guess for your partner, so choose your clue word wisely. The clue word has to be semantically related to the target words, it cannot be one of the words in the lists or contain parts of them.} \\
\\ 
\texttt{Always give your single{-}word clue and your comma{-}separated list of related target words in the following format and make your answers as short as possible, never include any other text than is required in this form:} \\
\\ 
\texttt{CLUE: <WORD>} \\
\texttt{TARGETS: <TARGETS>} \\
\\ 
\texttt{Your team words are: teacher, phoenix, plate.} \\
\texttt{Your opponent's team words are: root.} \\
\texttt{Distractor words are: heart, plastic, war, ninja, millionaire, snowman, orange.} \\
\texttt{The kill words are: .} \\
            }
        }
    }
    & \\ \\

    \theutterance \stepcounter{utterance}  
    & \multicolumn{4}{p{0.6\linewidth}}{
        \cellcolor[rgb]{0.9,0.9,0.9}{
            \makecell[{{p{\linewidth}}}]{
                \texttt{\tiny{[P1$\rangle$GM]}}
                \texttt{CLUE: education} \\
\texttt{TARGETS: teacher} \\
            }
        }
    }
    & & \\ \\

    \theutterance \stepcounter{utterance}  
    & & & \multicolumn{2}{p{0.3\linewidth}}{
        \cellcolor[rgb]{0.9,0.9,0.9}{
            \makecell[{{p{\linewidth}}}]{
                \texttt{\tiny{[GM$|$GM]}}
                \texttt{education} \\
            }
        }
    }
    & & \\ \\

    \theutterance \stepcounter{utterance}  
    & & & \multicolumn{2}{p{0.3\linewidth}}{
        \cellcolor[rgb]{0.9,0.9,0.9}{
            \makecell[{{p{\linewidth}}}]{
                \texttt{\tiny{[GM$|$GM]}}
                ['teacher']
            }
        }
    }
    & & \\ \\

    \theutterance \stepcounter{utterance}  
    & & & \multicolumn{2}{p{0.3\linewidth}}{
        \cellcolor[rgb]{0.9,0.9,0.9}{
            \makecell[{{p{\linewidth}}}]{
                \texttt{\tiny{[GM$|$GM]}}
                {'word': 'teacher', 'assignment': 'team'}
            }
        }
    }
    & & \\ \\

    \theutterance \stepcounter{utterance}  
    & & \multicolumn{4}{p{0.6\linewidth}}{
        \cellcolor[rgb]{0.9,0.9,0.9}{
            \makecell[{{p{\linewidth}}}]{
                \texttt{\tiny{[GM$\rangle$P2]}}
                \texttt{You have attempted this same task before and did not succeed. These are the plans you wrote after those attempts.} \\
\texttt{Trial \#1: New plan: I failed because I included the target word 'school' in my guess list, which is forbidden as the Game Master only provides a list of candidate words to choose from. Next time, I will strictly select only words from the provided list that semantically relate to the target concept, ensuring I never include the target word itself or any word not explicitly listed. For the current task ('school'), I will choose 'teacher' and 'ninja' (as a metaphorical or specific role association) or 'heart' (metaphorical), but strictly avoiding 'school'. Given the remaining list, 'teacher' is the strongest link. I will output: GUESS: teacher.} \\
\\ 
\\ 
\texttt{Provide a comma{-}separated list of up to 1 words from the following list that best relate or are most closely associated with the word 'education'. Always start your list of guess(es) with 'GUESS: ' and do not include any other text in your answer.} \\
\\ 
\texttt{heart, millionaire, phoenix, snowman, plate, plastic, root, teacher, ninja, war, orange} \\
            }
        }
    }
    & \\ \\

    \theutterance \stepcounter{utterance}  
    & & & \multicolumn{4}{p{0.6\linewidth}}{
        \cellcolor[rgb]{0.9,0.9,0.9}{
            \makecell[{{p{\linewidth}}}]{
                \texttt{\tiny{[GM$\langle$P2]}}
                \texttt{GUESS: teacher} \\
            }
        }
    }
     \\ \\

    \theutterance \stepcounter{utterance}  
    & & & \multicolumn{2}{p{0.3\linewidth}}{
        \cellcolor[rgb]{0.9,0.9,0.9}{
            \makecell[{{p{\linewidth}}}]{
                \texttt{\tiny{[GM$|$GM]}}
                ['teacher']
            }
        }
    }
    & & \\ \\

    \theutterance \stepcounter{utterance}  
    & & & \multicolumn{2}{p{0.3\linewidth}}{
        \cellcolor[rgb]{0.9,0.9,0.9}{
            \makecell[{{p{\linewidth}}}]{
                \texttt{\tiny{[GM$|$GM]}}
                {'word': 'teacher', 'assignment': 'team'}
            }
        }
    }
    & & \\ \\

    \theutterance \stepcounter{utterance}  
    & & & \multicolumn{2}{p{0.3\linewidth}}{
        \cellcolor[rgb]{0.9,0.9,0.9}{
            \makecell[{{p{\linewidth}}}]{
                \texttt{\tiny{[GM$|$GM]}}
                {'word': 'teacher', 'assignment': 'team'}
            }
        }
    }
    & & \\ \\

    \theutterance \stepcounter{utterance}  
    & & & \multicolumn{2}{p{0.3\linewidth}}{
        \cellcolor[rgb]{0.9,0.9,0.9}{
            \makecell[{{p{\linewidth}}}]{
                \texttt{\tiny{[GM$|$GM]}}
                {'word': 'root', 'assignment': 'opponent'}
            }
        }
    }
    & & \\ \\

    \theutterance \stepcounter{utterance}  
    & & & \multicolumn{2}{p{0.3\linewidth}}{
        \cellcolor[rgb]{0.9,0.9,0.9}{
            \makecell[{{p{\linewidth}}}]{
                \texttt{\tiny{[GM$|$GM]}}
                {'hidden': {'team': ['phoenix', 'plate'], 'innocent': ['heart', 'plastic', 'war', 'ninja', 'millionaire', 'snowman', 'orange'], 'opponent': [], 'assassin': []}, 'revealed': {'team': {'team': ['teacher'], 'innocent': [], 'opponent': [], 'assassin': []}, 'opponent': {'team': [], 'innocent': [], 'opponent': ['root'], 'assassin': []}}}
            }
        }
    }
    & & \\ \\

    \theutterance \stepcounter{utterance}  
    & & \multicolumn{4}{p{0.6\linewidth}}{
        \cellcolor[rgb]{0.9,0.9,0.9}{
            \makecell[{{p{\linewidth}}}]{
                \texttt{\tiny{[P1$\langle$GM]}}
                \texttt{The word teacher was guessed correctly. Your teammate's turn ended there.} \\
\texttt{Now provide another clue relating to some of your remaining team words and a list of the related target words. Remember to start your clue with 'CLUE: ', put a new line, and start your comma{-}separated list of target words with 'TARGETS: '. Notice: some words have been removed from the lists compared to previous requests.} \\
\\ 
\texttt{Your remaining team words are: phoenix, plate.} \\
\texttt{Remaining words for your opponent are: root.} \\
\texttt{Remaining distractor words are: heart, plastic, war, ninja, millionaire, snowman, orange.} \\
\texttt{Remaining kill words are: .} \\
            }
        }
    }
    & \\ \\

    \theutterance \stepcounter{utterance}  
    & \multicolumn{4}{p{0.6\linewidth}}{
        \cellcolor[rgb]{0.9,0.9,0.9}{
            \makecell[{{p{\linewidth}}}]{
                \texttt{\tiny{[P1$\rangle$GM]}}
                \texttt{CLUE: fire} \\
\texttt{TARGETS: phoenix} \\
            }
        }
    }
    & & \\ \\

    \theutterance \stepcounter{utterance}  
    & & & \multicolumn{2}{p{0.3\linewidth}}{
        \cellcolor[rgb]{0.9,0.9,0.9}{
            \makecell[{{p{\linewidth}}}]{
                \texttt{\tiny{[GM$|$GM]}}
                \texttt{fire} \\
            }
        }
    }
    & & \\ \\

    \theutterance \stepcounter{utterance}  
    & & & \multicolumn{2}{p{0.3\linewidth}}{
        \cellcolor[rgb]{0.9,0.9,0.9}{
            \makecell[{{p{\linewidth}}}]{
                \texttt{\tiny{[GM$|$GM]}}
                ['phoenix']
            }
        }
    }
    & & \\ \\

    \theutterance \stepcounter{utterance}  
    & & & \multicolumn{2}{p{0.3\linewidth}}{
        \cellcolor[rgb]{0.9,0.9,0.9}{
            \makecell[{{p{\linewidth}}}]{
                \texttt{\tiny{[GM$|$GM]}}
                {'word': 'phoenix', 'assignment': 'team'}
            }
        }
    }
    & & \\ \\

    \theutterance \stepcounter{utterance}  
    & & & \multicolumn{2}{p{0.3\linewidth}}{
        \cellcolor[rgb]{0.9,0.9,0.9}{
            \makecell[{{p{\linewidth}}}]{
                \texttt{\tiny{[GM$|$GM]}}
                \texttt{opponent has won} \\
            }
        }
    }
    & & \\ \\

    \theutterance \stepcounter{utterance}  
    & \multicolumn{4}{p{0.6\linewidth}}{
        \cellcolor[rgb]{0.95,0.95,0.95}{
            \makecell[{{p{\linewidth}}}]{
                \texttt{\tiny{[P2$\rightarrow$P1]}}
                \texttt{Post{-}episode Reflexion (not sent to GM):} \\
\texttt{New plan: I failed because I provided a clue ("fire") that only related to one remaining word ("phoenix") while ignoring the other ("plate"), allowing the opponent to guess my remaining word on their next turn. The strategy of targeting only one word per turn was too slow. Next time, I will select a single clue word that semantically connects to *all* of my remaining team words (e.g., "material" for "plate" and "phoenix" if a link exists, or finding a broader concept) to clear multiple targets in one move, ensuring I win before the opponent can guess the rest.} \\
            }
        }
    }
    & & \\ \\

    \theutterance \stepcounter{utterance}  
    & & & \multicolumn{4}{p{0.6\linewidth}}{
        \cellcolor[rgb]{0.95,0.95,0.95}{
            \makecell[{{p{\linewidth}}}]{
                \texttt{\tiny{[P1$\rightarrow$P2]}}
                \texttt{Post{-}episode Reflexion (not sent to GM):} \\
\texttt{New plan: I failed because I only provided a single word ('teacher') when the task allowed for up to one word, but I likely missed the opportunity to include the most direct synonym available in the list, or I misinterpreted the constraint regarding the count. Upon re{-}evaluating the list, 'root' is a strong metaphorical link (root of education), but 'teacher' is the most direct association. However, looking closely at the previous failure context where I might have included forbidden words or missed the optimal single choice: The list contains 'teacher', which is the best fit. If the previous attempt failed despite being 'teacher', it might be because the system expected a specific format or I hallucinated a second word in a different turn. Wait, the prompt says "up to 1 words". I provided one. The failure reason "failed to achieve the target state" implies 'teacher' was not the correct answer or the only correct answer expected. Let's re{-}examine the list for a better semantic match. 'root' (foundation), 'heart' (core), 'war' (no), 'ninja' (no). Actually, 'teacher' is the standard association. Perhaps the target word was something else entirely? No, the task is to relate to 'education'. Maybe the issue is that I need to be more confident in the single best choice. Let's reconsider the list: 'root' could be the "root" of education. But 'teacher' is the agent. If 'teacher' failed, maybe the intended answer is 'root' as in the foundation? Or perhaps I need to check if 'school' was in the list? No, 'school' is not in the provided list. The list is: heart, millionaire, phoenix, snowman, plate, plastic, root, teacher, ninja, war, orange. 'Teacher' is the only logical direct link. If 'teacher' failed, maybe the game master considers 'root' (as in the root cause or foundation) or 'heart' (education of the heart) as the intended answer? Or perhaps I need to output the word that is *most* closely associated, and 'teacher' was rejected for some reason? Let's assume the target association is 'root' (foundation) or 'heart'. Given the ambiguity, I will select 'root' as it represents the foundation of education, which is a very strong conceptual link, whereas 'teacher' is a role. Wait, 'teacher' is usually the best. Let's look at the failure again. "failed to achieve the target state". This means my guess was wrong. If 'teacher' is wrong, what is the next best? 'root' (the root of education). I will switch to 'root'.} \\
\\ 
\texttt{Actually, looking at the list again, 'root' is often used metaphorically ("root of the problem", "root of education"). 'Teacher' is a person.} \\
            }
        }
    }
     \\ \\

\end{supertabular}
}

\end{document}
